% Contexto e instrucciones literales de EXA-0007, páginas 1 y 2.
% Fuente: Selectividad/matematicas_aplicadas/ORD_TITULAR_MATEMÁTICAS APLICADAS CCSS-JUNIO-IMPRESO.pdf
% SHA-256: c42a605501a277fe7194b964ad92ad66047be9a62dc254fa2219d756d3d71317
\section*{Contexto del examen original}
\textbf{PRUEBA DE ACCESO Y ADMISIÓN A LA UNIVERSIDAD}\par
\textbf{ANDALUCÍA, CEUTA, MELILLA y CENTROS en MARRUECOS}\par
CONVOCATORIA ORDINARIA, CURSO 2020-2021\par
\textbf{MATEMÁTICAS APLICADAS A LAS CIENCIAS SOCIALES II}

\textbf{Instrucciones:}
\begin{enumerate}[label=\textbf{\alph*)}]
  \item Duración: 1 hora y 30 minutos.
  \item Elija \underline{cuatro} de los ocho ejercicios propuestos de
  \underline{al menos tres bloques distintos}. Se corregirán los cuatro
  primeros ejercicios que aparezcan en el examen y que cumplan el requisito
  anterior.
  \item En cada ejercicio, parte o apartado se indica la puntuación máxima
  asignada.
  \item Todos los resultados deben estar suficientemente justificados.
  \item Se permitirá el uso de calculadoras que no sean programables,
  gráficas ni con capacidad para almacenar o transmitir datos. Si obtiene
  resultados directamente con la calculadora, explique con detalle los pasos
  necesarios para su obtención sin el uso de la misma.
\end{enumerate}
\textbf{\underline{Este examen consta de 4 Bloques (A, B, C y D)}}\par
Deberá responder a cuatro ejercicios de entre los ocho propuestos con la
condición de que pertenezcan al menos a 3 bloques distintos. En caso de
responder a más ejercicios de los requeridos, serán tenidos en cuenta los
respondidos en primer lugar.

\noindent\textbf{Nota de la colección:} Se reúnen los ocho ejercicios para
consulta. La regla anterior corresponde al examen original. La convocatoria
ordinaria y el curso están impresos; el tipo titular y junio sólo proceden
del nombre del archivo. El nombre del auxiliar de criterios dice Modelo-2,
pero ninguna de las cabeceras revisadas imprime un número de modelo.
\par\medskip
